\section{Average Linkage HAC on Trees is \texorpdfstring{\CC}{CC}-Complete}
In this section, we show that Average Linkage HAC on trees is in \CC, and hence, combined with the \CC-Hardness result, proves that it is \CC-Complete. In particular, we show that Average Linkage HAC on trees can be (log-space) reduced to a polynomial-many instances of a generalization of Adaptive Minimum, which we also show is \CC-Complete. We first recall a more general matching problem that is known to be \CC-Complete.

\begin{problem}[\lfbm]\label{dfn:LFBMM} An instance of lexicographically first maximal b-matching (LFM B-Matching) consists of a graph $G=(V, E)$ with vertices ordered as $V = (v_0, v_1, \ldots, v_{n-1})$. Further, for each vertex $v$ we are given a number $b(v)$. The lexicographically first maximal b-matching is the matching obtained by running the greedy algorithm for maximal matching on edges ordered first by their lower endpoint and then their higher endpoint such that the number of edges in the matching incident on any vertex $v$ doesn't exceed $b(v)$. That is, in this ordering $e = \{v_i, v_j\}$, with $i<j$, precedes $e' = \{v_{i'}, v_{j'}\}$, with $i'<j'$, iff (1) $i < i'$ or (2) $i = i'$ and $j < j'$. If any ties persist, i.e. multi-edges exist, break them arbitrarily. Our goal is to decide if a designated input edge is in the LFM B-Matching.
\end{problem}
In the case of multi-edges, any copy being in the matching is sufficient. Alternatively, we can also have a pre-defined ordering on the copies and break ties according to this ordering. 

\begin{theorem}\label{thm:LFBMCCComplete}
    \lfbm is \CC-Complete.
\end{theorem}
The proof follows by a standard reduction of B-Matching to Matching, and the fact that LFM Matching on general graphs is \CC-Complete~\cite{mayr1992complexity}.

\begin{problem}[\amj]\label{l:amj}
    An instance of \amj consists of a ($0$-based indexed) $n \times n$ matrix $A$ where each cell consists of two values: weight ($A[i,j].\weight$), and jump ($A[i,j].\jump$). For each row, the weights form a permutation of $\{0, \ldots, n-1\}$, whereas the jump values lie in the range $[0,n)$. We are also given some index $x \in [0, n)$. The goal is to simulate the following algorithm.
    Start with $I = \{0, \ldots, n-1\}$ and $i=0$, execute the following steps while $i \le x$:
    \begin{enumerate}
        \item Let $k_i = \argmin_{j \in I} A[i, j].\weight$.
        \item Set $I := I \setminus \{k_i\}$.
        \item For each $\ell=i+1$ to $i+A[i,k_i].\jump$, set $k_\ell = \Null$.
        \item Set $i = i + A[i, k_i].\jump$.
    \end{enumerate}
    Our goal is to compute $k_x$.
\end{problem}

\begin{lemma}
    \amj is \CC-Complete.
\end{lemma}
\begin{proof}
    Observe that \am is a special case of \amj when all the jump values are $1$. Thus, \amj is also \CC-Hard. We now prove that it is in \CC by reducing it to the \lfbm problem. 
    
    For each row $i \in [0,n)$ we introduce a node $r_i$. For each pair $(i,j) \in [0,n)\times[0,n)$, we introduce a column node $c_{i,j}$ and a corresponding auxiliary node $a_{i,j}$. Let $k$ be a parameter whose value will be determined later. We set the capacities $b(r_i) = b(c_{i,j}) = b(a_{i,j}) = k/{(i+1)^2}$. Next, we introduce 4 types of edges: 
    \begin{enumerate}[label=(\roman*)]
        \item Between nodes $r_i$ and $c_{i,j}$, introduce $k/(i+1)^2$ edges.
        \item Between nodes $c_{i,j}$ and $a_{i,j}$, introduce $k/(i+1)^2$ edges.
        \item Between nodes $a_{i,j}$ and $c_{i+1,j}$, introduce $k/(i+2)^2$ edges.
        \item Let $t = A[i,j].\jump$. If $t>0$, for each $\ell=1$ to $t-1$, introduce $k/(i+\ell+1)^2$ edges between nodes $a_{i,j}$ and $r_{i+\ell}$.
    \end{enumerate}
    Finally, we define the ordering on the vertices. For row $i$, the first node in the ordering will be $r_i$, followed by the column nodes $c_{i,j}$, for $j\in [0,n)$, ordered by their weights in $A[i,j]$. Next, the $n$ auxiliary nodes $a_{i,j}$ follow, again ordered by their corresponding weights in $A[i,j]$. The ordering of row $i+1$ is placed after the ordering of row $i$, thus obtaining the overall ordering on nodes.

    Intuitively, the auxiliary node acts as a switch that is triggered when its corresponding column node doesn't match with it, i.e. the column node matches with the row node. When it is triggered, it blocks all future column nodes from participating in matches with row nodes via a ``chain-reaction''. It also matches and saturates the capacities of all row nodes that get skipped due to the jump value at that cell, thus simulating a jump.

    We claim that $k_i = j$ iff $r_i$ is matched to $c_{i,j}$ ($b(r_i)$ times), and $k_i = \Null$ iff $r_i$ is matched to $a_{i',k_{i'}}$ for some $i'<i$ ($b(r_i)$ times). We prove this by induction on $i$. Consider the index $i$. Let $i'$ denote the previous index that was matched to a column node $c_{i',j'}$. Let $t = A[i',j'].\jump$. By induction, we know that $k_{i'}=j'$, and $k_{s} = \Null$ for all $s \in [i'+1,i-1]$ implying that $t \ge i-i'-1$. If $t \ge i-i'$, then observe that $k_i=\Null$. Here, 



\end{proof}

\begin{theorem}\label{thm:avgtreeCCComplete}
    Average linkage HAC on trees is \CC-Complete.
\end{theorem}
\begin{proof}
    By \Cref{thm:mainLowerPar}, we know that it is \CC-hard. We now prove it is in \CC. We follow a similar framework used by \Cref{alg:paths} for the Paths case. Bucket the edges by weights and process them one-by-one. Consider the subgraph obtained by the graph induced on the edges in the current bucket. If its a forest with multiple connected components, we can process each component independently. So, we can assume that we have a tree. Root the tree at the largest vertex by size. Now, we decompose the tree into \emph{blocks} as follows: If the size of a vertex is at most $1/2$ the size of its parent, then include it in its parent's block. Otherwise, start a new block at that vertex. Observe that each vertex belongs to exactly one block, and since the cluster sizes reduces by a constant factor as we go down, the height of the subtree induced on any block will be at most $\log n$. If we contract each block into a single vertex, we get a tree representation of this block decomposition. Note that the structure of this tree can be arbitrary. For simplifying the presentation, we start with the assumption that we have only one block. We remove these assumptions gradually.

    Consider the case when the block has height $2$, i.e., we have a star graph rooted at the center. How do we process a star? Since these edges share an endpoint, they will be merged sequentially. Further, they will be merged in the sorted order of their initial weights since only the common endpoint's size is changing, and thus, the relative ordering of these edges will be preserved.

    Now, consider a block of height $3$. Here, we have a root, its children whom we call internal nodes, and their children whom we call leaves.
    % The subtree rooted at an internal node behaves similar to the star case discussed above, except that sometimes the internal node will prefer to merge with the root. If the root prefers that child at the same time, we can merge them and the remaining leaves in that subtree can be dropped as the edges from the root to these leaves will have its weight dropped by a constant factor after this merge.

    Now we will consider a path of blocks rooted at one of the end blocks. Look at one of the inter-block edges $(u,v)$. If $(u,v)$ decides to merge with the block of $u$, then the weight of every edge incident to $v$ in the block of $v$ drops by a constant factor (similar to the cases obtained in \Cref{alg:paths}). Thus, starting with the first block, we can propagate the merging behavior of these inter-block edges in a similar manner as in \Cref{lem:path_prefix_sum} using a prefix-sum routine. 

    Finally, we return to the original arbitrary tree structure on the blocks. Given this, we compute a heavy-light decomposition (solvable in NC) on the tree structure. We start by applying the path algorithm discussed above to the leaf-paths.

\end{proof}